\begin{abstract}
我们设计并实现了Google文件系统（GFS），这是一个面向大规模
分布式、数据密集型应用的可扩展分布式文件系统。它运行在廉价
的通用硬件上，同时提供容错能力，并能为大量客户端提供很高的
聚合性能。

虽然与之前的分布式文件系统有许多相同的目标，但是我们的设计主要受到两方面观察的驱动：
一是我们应用的实际工作负载/访问模式，二是当时及预期中的技术环境。这与之前的一
些文件系统假设有显著的不同。因此，我们重新审视了传统选择，并探索了截然不同的设计方向。

该文件系统已经成功满足了我们的存储需求。
它已在Google内部被广泛部署，作为一种存储平台，用于生成和处理我们
服务所需的数据，也用于需要大规模数据集的研究与开发工作。
到目前为止，最大的集群在一千多台机器上的数千块磁盘中提供了数百TB的存
储容量，并且有数百个客户端同时访问它。

本论文将介绍为支持分布式应用而扩展的文件系统接口，讨论设计中的
多个方面，并报告微基准测试与真实使用场景下的测量结果。
\end{abstract}
